\section{Related Work}
\label{sec:related}

\paragraph{Continual learning and catastrophic forgetting.}
Catastrophic forgetting---the tendency of neural networks to lose previously learned knowledge when trained on new tasks---has been studied extensively since \citet{mccloskey1989catastrophic} and \citet{french1999catastrophic}.
The continual learning literature broadly organizes mitigations into three families: regularization-based, replay-based, and architecture-based methods~\citep{delange2021continual}.
SFP draws from all three: it regularizes activations (not weights), uses a small memory buffer for basis construction, and operates within a fixed architecture.

\paragraph{Regularization methods.}
Elastic Weight Consolidation~\citep[EWC;][]{kirkpatrick2017overcoming} penalizes changes to parameters deemed important for previous tasks, using the diagonal of the Fisher information matrix as an importance measure.
Synaptic Intelligence~\citep{zenke2017continual} computes importance online from the path integral of gradients.
These methods operate in \emph{parameter space}, regularizing weight changes.
SFP instead operates in \emph{activation space}, which is lower-dimensional and more directly tied to model behavior.
Moreover, EWC's diagonal Fisher approximation can be poor for modern overparameterized models, while SFP's PCA basis adapts to the actual activation geometry.

\paragraph{Knowledge distillation.}
Learning without Forgetting~\citep[LwF;][]{li2017learning} distills the old model's output logits into the new model during fine-tuning.
PODNet~\citep{douillard2020podnet} extends this to intermediate feature maps, penalizing changes in pooled hidden-state statistics.
SFP is related to PODNet in that both preserve intermediate representations, but SFP is \emph{selective}: rather than penalizing all representational change, it identifies and preserves only the dimensions that matter for old-task performance, leaving the model free to learn in the complementary subspace.

\paragraph{Replay and memory.}
Experience Replay~\citep{chaudhry2019tiny} maintains a small buffer of old-task examples and interleaves them during training.
A-GEM~\citep{chaudhry2018efficient} uses the memory buffer to project gradients away from directions that would increase old-task loss.
SFP uses a memory buffer, but only for \emph{basis construction} (extracting PCA directions and computing ablation importance), not for direct training.
This distinction means SFP's memory requirements are small and fixed: the buffer is used once to build the basis, then discarded during training.

\paragraph{Architectural methods.}
Progressive Neural Networks~\citep{rusu2016progressive} allocate new parameters for each task.
Orthogonal LoRA methods~\citep{qin2024exploring} constrain low-rank adapter updates to be orthogonal to the subspace of previous adapters, preventing interference by construction.
SFP shares the geometric intuition of orthogonal methods---learning should occur in a subspace complementary to important features---but identifies the important subspace empirically via ablation rather than imposing it through adapter orthogonality.
Our \texttt{gradient\_in\_complement} theorem (\cref{sec:formal}) formalizes this connection: at a preserved stationary point, the new-task gradient lies entirely in the complement of the preserved subspace.

\paragraph{Mechanistic interpretability.}
Recent work on sparse autoencoders~\citep{bricken2023monosemanticity,cunningham2023sparse} and activation patching~\citep{meng2022locating,wang2023interpretability} has shown that individual directions in activation space can encode interpretable, causally relevant features.
SFP's ablation-based importance ranking is conceptually related to activation patching: both assess the causal role of activation directions by measuring downstream effects of interventions.
The connection to superposition theory~\citep{elhage2022toy} motivates our hypothesis that forgetting may be low-dimensional even in models that represent many features in superposition.

\paragraph{Formal verification in machine learning.}
Formal verification of neural network properties has focused primarily on robustness certification~\citep{katz2017reluplex,wong2018provable}.
To our knowledge, SFP is the first continual learning method with formally verified mathematical properties.
Our Lean~4 formalization~\citep{mathlib2020} verifies the complete reduction from the empirical hypothesis (H1) to a forgetting bound, making the logical structure of the argument machine-checkable and the assumptions explicit.
